\documentclass[10pt,a4paper]{amsart}
\usepackage[margin=19mm]{geometry}
\usepackage{amsmath,amssymb,mathtools}
\usepackage[scr=rsfs,cal=euler]{mathalfa}
\usepackage{xfrac}
\usepackage[unicode,hidelinks,bookmarks=false]{hyperref}
\newcommand{\ie}{i.e.\ }
\newcommand{\cf}{cf.\ }
\newcommand{\resp}{respectively\ }
\newcommand{\Subsection}{\subsection}
\providecommand{\nobreakdash}{\nobreak\hbox{-}}
\setlength{\emergencystretch}{2em}
\multlinegap0pt
\usepackage{tikz}
\usepackage{float}
\newtheorem{lemma}{Lemma}
\newcommand{\eqn}[1]{\begin{equation*} #1 \end{equation*}}
\newcommand{\neqn}[1]{\begin{equation} #1 \end{equation}}
\title{Minimal Bridges and a Rotation-Based Bijection}
\author{Benjamin Lou, Lucas Augustus Brown}
\date{Source: arXiv:2608.11005v1 (CC BY 4.0)}
\begin{document}
\maketitle

\noindent\emph{Source and licence.} Minimal Bridges and a Rotation-Based Bijection. Version arXiv:2608.11005v1 (CC BY 4.0), distributed by arXiv under the Creative Commons Attribution 4.0 International licence, http://creativecommons.org/licenses/by/4.0/, as recorded in the arXiv licence metadata for this exact version and shown on the abstract page https://arxiv.org/abs/2608.11005v1. Full licence notice: \texttt{licenses/arxiv-2608.11005-CC-BY-4.0.txt}.\par\smallskip

\noindent\textit{Editorial scope.} Counted unit: the article's lemma lemma:KEOO (source0.tex lines 494--496). The article's complete original proof of it is delivered with it (source0.tex lines 497--507, one \texttt{proof} environment of 11 lines). No mathematics is written, repaired or re-derived: the statement and the proof are the article's own lines, in the article's own notation, and no retained line is substituted or re-flowed.  Retained uncounted as the source-local context the proof needs: lines 483--485.  Nothing was omitted from any retained span, so the package carries no omission record.  No retained line contains a citation, so the wrapper carries no bibliography and no bibliography entry.  No retained line contains a figure environment, an image-inclusion command or drawing code, so no image is delivered and the package declares no asset dependency.  Editorial formatting only, in addition: an AMS \texttt{amsart} preamble that restates the article's own declarations and macros used by the retained lines, namely the environments \texttt{lemma} on separate counters, together with the standard packages those lines need.  The retained environments print in this document as Lemma 1 in order of appearance, whereas the article prints the same items with the numbering of its own full document.  The complete native source of the article as distributed in the arXiv e-print for this exact version is bundled unchanged as \texttt{sources/207378/source0.tex}.
\begin{lemma} \label{lemma:KnullDO}
$K(\emptyset,D,O) = 2C_{2n}$.
\end{lemma}

\begin{lemma} \label{lemma:KEOO}
$K(E,O,O) = C_{2n} + 2C_{2n-1}$.
\end{lemma}

\begin{proof}
The paths in $\mathcal{K}(\emptyset, D, O)$ can be partitioned into two subsets based on whether their first diagonal crossing is an even or an odd lattice point; therefore,
\eqn{K(\emptyset,D,O) = K(E,O,O) + K(O,E,O).}
The LHS is $2C_{2n}$ by Lemma \ref{lemma:KnullDO}. Therefore, it suffices to show that
\neqn{K(O,E,O) = C_{2n} - 2C_{2n-1}. \label{opaujnertfb}}
The set $\mathcal{K}(O,E,O)$ is the same as $\mathcal{K}(\emptyset,\emptyset,O)$ but without those paths that do not touch the diagonal. Therefore,
\eqn{K(O,E,O) = K(\emptyset,\emptyset,O) - K(\emptyset,\emptyset,D).}
Now $\mathcal{K}(\emptyset,\emptyset,O)$ is the set of Shapiro paths. Meanwhile, $\mathcal{K}(\emptyset,\emptyset,D)$ is the set of paths that stay strictly above or below the diagonal, so\footnote{The paths strictly below the diagonal are the Dyck paths with $2n-1$ rightward steps, plus an initial rightward step and a final upward step, and the paths above the diagonal are the half-turn rotations thereof.}
\neqn{\mathcal{K}(\emptyset,\emptyset,D)= 2C_{2n-1}. \label{non-diag}}
This proves \ref{opaujnertfb} and hence the result.
\end{proof}

\end{document}
